\documentclass[12pt]{article}
\usepackage{amsmath, amssymb, amsthm}
\usepackage{graphicx}
\title{Discrete Math HW2}
\author{Anatolii Larkin}

\begin{document}

\maketitle


\begin{enumerate}

    \item \textit{Geometric progression} --- a sequence of non-zero numbers $b_1, b_2, b_3, \dots$, where each term is found by multiplying the previous one by constant number $q$: $b_{i+1} = b_i \cdot q$. To define a sequence it is sufficient to set the first element $b_1$ and the common ratio $q$.
    
    Let $S_n = b_1 + b_2 + b_3 + \dots + b_n$, where $n \in \mathbb{N}$ and $b_{i+1} = b_i \cdot q$ with $q \neq 1$.
    
    \medskip
    We can express each term in terms of the first element $b_1$ and the common ratio $q$:
    \[
    S_n = b_1 + b_1q + b_1q^2 + \dots + b_1q^{n-1} = b_1(q^0 + q^1 + q^2 + \dots + q^{n-1})
    \]
    
    To make the main statement hold, we will first prove by mathematical induction that for all $n \in \mathbb{N}$:
    \[
    q^0 + q^1 + q^2 + \dots + q^{n-1} = \frac{1 - q^n}{1 - q}
    \]
    
    \noindent{Base Case ($n = 1$):} \\
    The left-hand side is $q^0 = 1$. The right-hand side is:
    \[
    \frac{1 - q^1}{1 - q} = 1
    \]
    Thus, the base case holds.
    
    \medskip
    \noindent{Inductive Step:} \\
    Assume that the statement holds for a specific integer $k \ge 1$:
    \[
    q^0 + q^1 + \dots + q^{k-1} = \frac{1 - q^k}{1 - q}
    \]
    
    Now, we check the statement for $k + 1$:
    \begin{align*}
    q^0 + q^1 + \dots + q^{k-1} + q^k &= \frac{1 - q^k}{1 - q} + q^k \\[1ex]
    &= \frac{(1 - q^k) + q^k(1 - q)}{1 - q} \\[1ex]
    &= \frac{1 - q^k + q^k - q^{k+1}}{1 - q} \\[1ex]
    &= \frac{1 - q^{k+1}}{1 - q}
    \end{align*}
    Therefore, the inductive hypothesis holds for $k + 1$. By the principle of mathematical induction, the identity is true for all $n \in \mathbb{N}$.
    
    \medskip
    Substituting this identity back into our equation for $S_n$:
    \[
    S_n = b_1 \left( \frac{1 - q^n}{1 - q} \right)
    \]
    Then the statement holds. 



    \item 

    Prove that for every natural number \(n\) the equation holds true:
    \[
    1^2 + 2^2 + \dots + n^2 = \frac{n(n+1)(2n+1)}{6}.
    \]


    Consider base case n = 1.

    Left part = $1^2 = 1$
    
    Right part = $\frac{1*2*3}{6} = 1$
    
    Base case holds

    Inductive hypothesis: suppose that for $k \geq 1$ ; $k \in N$, the following statement holds:

    1^2 + 2^2 + ... + k^2 = \frac{k(k+1)(2k+1)}{6}

    Then for k + 1:

    1^2 + 2^2 + ... + k^2 + (k+1)^2 = \frac{k(k+1)(2k+1)}{6} + (k+1)^2 = \frac{(k+1)(2k^2 + 7k + 6)}{6} = \frac{(k+1)(k+2)(2k+3)}{6} = \frac{(k+1)((k+1)+1)(2(k+1)+1)}{6}

    Therefore the statement holds


    \item 

    Count the number of integers between 1 and 1000, that are divisible by 6 and 10

    Consider the set A = \{n \| $n \in N$ and $n \in (1,1000)$ and $n = 6k$, $k \in N$ and $n = 10p$, $p\in N$\}

    Since divisors 10 and 6 are not prime, perform prime number factorization to get LCM (6,10):

    6 = 2 * 3
    
    10 = 2 * 5 

    LCM(6,10) = 2 * 5 * 3 = 30

    Then A = \{n \| $n \in N$ and $n \in (1,1000)$ and $n = 6k$, $k \in N$ and $n = 10p$, $p\in N$\} = \{n \| $n \in N$ and $n \in (1,1000)$ and $n = 30q$, $q \in N$ \}

    Every 30th number in sequence (1,1000) of natural numbers is divisible by 30.

    Then \|A\| = $\lfloor\frac{1000}{30}\rfloor$ = 33


    \item  Count the number of integers between 1 and 1000, that are divisible by either
3 or 10, but are not divisible by 15.

 Consider sets:
    A &= \{n \mid n \in \mathbb{N} \text{ and } 1 \le n \le 1000 \text{ and } n = 3k, \, k \in \mathbb{N}\} \\
    B &= \{n \mid n \in \mathbb{N} \text{ and } 1 \le n \le 1000 \text{ and } n = 10p, \, p \in \mathbb{N}\} \\
    C &= \{n \mid n \in \mathbb{N} \text{ and } 1 \le n \le 1000 \text{ and } n = 15q, \, q \in \mathbb{N}\}

By computing the sizes of these sets:

    \|A\| &= \left\lfloor\frac{1000}{3}\right\rfloor = 333 \\
    \|B\| &= \left\lfloor\frac{1000}{10}\right\rfloor = 100 \\
    \|C\| &= \left\lfloor\frac{1000}{15}\right\rfloor = 66

Number that are divisible either by 3 or 10, but not 15 belong to set $X = (A \cup B) \setminus C$.

By the inclusion-exclusion principle for set subtraction:
\[

\|X\| = \|A\| + \|B\| - \|A \cap B\| - \|A \cap C\| - \|B \cap C\| + \|A \cap B \cap C\|
\]

Where the intersections are defined as:

    \|A \cap B\| &= \left\lfloor\frac{1000}{\text{lcm}(3,10)}\right\rfloor = \left\lfloor\frac{1000}{30}\right\rfloor = 33 \\
    \|A \cap C\| &= \left\lfloor\frac{1000}{\text{lcm}(3,15)}\right\rfloor = \left\lfloor\frac{1000}{15}\right\rfloor = 66 \\
    \|B \cap C\| &= \left\lfloor\frac{1000}{\text{lcm}(10,15)}\right\rfloor = \left\lfloor\frac{1000}{30}\right\rfloor = 33 \\
    \|A \cap B \cap C\| &= \left\lfloor\frac{1000}{\text{lcm}(3,10,15)}\right\rfloor = \left\lfloor\frac{1000}{30}\right\rfloor = 33


Substituting the values into the formula:

    \|X\| &= 333 + 100 - 33 - 66 - 33 + 33 \\
    \|X\| &= 334

\item Count the number of integers between 100 and 999, that have at least one
digit divisible by 4 (0 is divisible by any integer).

Suppose U is set of all possible numbers between 100-999 (900 numbers total). Suppose set A contains numbers with at least 1 digit divisible by 4. Suppose set B contains number with no digits divisible by 4.

Since no number can belong to A and B simultaneously, $|A| = |U| - |B|

Each possible digit for $x \in B$ can be found in set $\{1,2,3,5,6,7,9\}$, 7 digits total. Then |B| = 7^3 = 343

Then $|A| = 900 - 343 = 557$

\item Prove de Morgan’s laws:

\overline{A \cap B} &= \overline{A} \cup \overline{B} \\
\overline{A \cup B} &= \overline{A} \cap \overline{B}

Using direct proof and logical de Morgan's laws:

$(A \cap B)^c = \{x| x \not\in (A \cap B)\} = \{x | x \not\in A \lor x \not\in B\} = \{x | x \in A^c \lor x \in B^c\} = A^c \cup B^c$. Statement 1 is proven


$(A \cup B)^c = \{x | x \not\in (A \cup B)\} = \{x| x \not\in A \land x \not\in B\} = \{x| x \in A^c \land x \in B^c\} = A^c \cap B^c$. Statement 2 is proven

\item Simplify the following expression using set operations and de Morgan’s laws
in such a way, that the complement is only applied to A, B or C, but not to
any complex expression. Simplify the experssion.

$\neg(A\cup(\neg B \cap A)$ \backslash $(\neg C \cup \neg(B \cup A)))$

Left from backshlash:

(A \cup \neg B) \cap A = A

Right from backslash:

(\neg C \cup \neg(B \cup A)) = \neg C \cup (\neg B \cap \neg A)


Then $\neg(A\cup(\neg B \cap A)$ \backslash $(\neg C \cup \neg(B \cup A)))$ = \neg(A \backslash (\neg C \cup (\neg B \cap \neg A))) = \neg((A \cap C) \cup (A \cap B \cap C)) = \neg (A \cap C) = \neg A \cup \neg C

\item Draw the Venn diagram for the set $\neg A \cup \neg B \cup C$

\begin{figure}[h]
    \centering
    \includegraphics[width=0.75\linewidth]{2_8.png}
\end{figure}

The requested set is \textbf{ALL EXCEPT} orange compartment



\item For each of the following binary relations on the set \(A=\{a,b,c\}\) check, whether it is reflexive, anti-reflexive, symmetric, anti-symmetric, transitive:

1. \(R_1=\{(a,a),(a,b),(b,a),(b,b),(c,c)\}\)

Reflexive, symmetric,transitive 

2.  \(R_2=\{(a,a),(a,c),(c,b),(a,b),(b,b),(c,c)\}\)

Reflexive, anti symmetric, transitive

3.  \(R_3=\{(a,a),(a,c),(c,b),(a,b)\}\)

Anti-symmetric, transitive


\item Consider two sequences of sets \(A_i\) и \(B_i\), where \(i\) ranges from \(1\) to \(n\). Define
\(X = \bigcup_{i=1}^n A_i = A_1 \cup A_2 \cup \ldots \cup A_{n-1} \cup A_n\), \(Y = \bigcup_{i=1}^n B_i\), \(Z = \bigcup_{i=1}^n (A_i \setminus B_i)\). Come up with an example where \(X \setminus Y \neq Z\). Is any of inclusions \(X \setminus Y \subset Z\) or \(Z \subset X \setminus Y\) true?
\textit{P.S. To prove falsity, it is enough to provide a counterexample. To prove truth, a formal proof is required.}

1. Come with example where X \backslash Y \not = Z

Suppose n = 2; $A_1 = \{x\}; A_2 = \varnothing; B_1 = \varnothing; B_2 = \{x\}$

Then X = \{x\} \cup \varnothing = \{x\}

Y = \varnothing \cup \{x\} = \{x\}

Z = ($\{x\} \backslash \varnothing) \cup (\varnothing \backslash \{x\}$) = \{x\}

$X \backslash Y = \{x\} \backslash \{x\} = \varnothing \not = Z \longrightarrow$ counterexample found


2. Is inclusion $X \backslash Y \subset$ Z true?

Let y \in $X \backslash Y$. Then $(y \in X) \land (y \not \in Y)$. Then $(y \in A_p, p \in [1,n]) \land (y \not \in B_p, p \in [1,n])$

Since y can belong to any $A_p$, and y doesnt belong to any $B_p$, $y \in A_p \backslash B_p$.

Since Z is union of $A_i \backslash B_i, i \in [1,n]$, $y \in Z$. Then $\all y \in X \backslash Y$ also belong to Z, therefore inclusion $X \backslash Y \subset$ Z is true

3. Is inclusion $Z \subset X \backslash Y$ true?

This inclusion is not always true, see paragraph 1 of this task


    
    


\end{enumerate}

\end{document}